% Resumen en castellano (200-500 palabras según la guía de la ETSINF).
% Es el texto de tu main.tex de septiembre: hay que rehacerlo al final con
% % los resultados del bloque 1, el bloque 2 y el modelo final.
Resumen
% Las interfaces cerebro-computador para el habla buscan restaurar la comunicación en personas
% que han perdido la capacidad de hablar. Este trabajo estudia la viabilidad de adaptar OWSM,
% un modelo de reconocimiento de voz preentrenado de propósito general, a la decodificación de
% actividad neuronal intracortical del reto \textit{Brain-to-Text '25}. Se sustituye el
% \textit{frontend} acústico por un módulo adaptador entrenado desde cero, se incorpora una capa
% específica de sesión y se entrena el conjunto con un objetivo CTC puro, evaluando mediante
% ablaciones controladas qué componentes del preentrenamiento transfieren a una modalidad no
% acústica. El resultado principal es negativo y reproducible: un codificador inicializado
% aleatoriamente supera al inicializado con los pesos de OWSM (tasa de error por carácter de
% 0,312 frente a 0,381) con la mitad de coste, y un codificador aleatorio congelado iguala a uno
% preentrenado congelado. Dos hallazgos colaterales resultan de mayor alcance práctico: la
% granularidad de la unidad de salida produce un efecto cuatro veces superior al del
% preentrenamiento, favoreciendo unidades lingüísticamente gruesas frente a las articulatorias,
% y la degradación ante días de registro no vistos alcanza los 27 puntos porcentuales, con la
% capa de sesión agravándola. El sistema alcanza una tasa de error por carácter de 0,267 en
% condiciones comparables a la literatura reciente sobre el mismo conjunto de datos.
